% ===== 原文 pp.39--44 =====
% 说明：原文 p.43 为整页插图（图 10、图 11），“障碍类”方框被它夹在 p.42 与 p.44 之间；
%       按 STYLE_GUIDE §7 合并为一个方框，两幅插图置于方框之后。
% [原文 p.39]
\subsection{非阿贝尔规范理论}

在本小节中，我们使用基于屏蔽的论证来确定具有非阿贝尔规范群的规范理论的高阶形式对称性群。

\subsubsection{纯 $\SU(2)$ Yang--Mills 理论}

我们首先考虑一个 $d$ 维规范理论，其规范群为
\begin{equation}
  \cG=SU(2)\,, \tag{3.79}
\end{equation}
且不含物质场。

\para{Wilson 线算符。} Wilson 线算符由规范群 (3.79) 的不可约表示参数化：
\begin{equation}
  \cW_j=\Tr_{R_j}\mathcal{P}\ee^{2\pi i\int A}\,, \tag{3.80}
\end{equation}
其中 $\Tr_{R_j}$ 表示在自旋-$j$ 不可约表示 $R_j$ 中取迹，这里 $j$ 是一个半整数
\begin{equation}
  j\in\frac{1}{2}\Z\,. \tag{3.81}
\end{equation}

\para{利用场强构造非真局域算符。} 回忆场强 $F$ 在规范群的伴随表示中变换，而伴随表示就是自旋-1 表示
\begin{equation}
  j=1\,. \tag{3.82}
\end{equation}
% [原文 p.40]
这意味着场强的插入 $F(x)$ 不能给出一个真局域算符，因为它将不是规范不变的。取而代之，我们可以把 $F(x)$ 附着到伴随表示中的 Wilson 线 $\cW_1$ 上，得到一个规范不变的构型
\begin{equation}
  F(x)\cW_1(L)\,, \tag{3.83}
\end{equation}
其中 $L$ 是一条终止于点 $x$ 的线，即
\begin{equation}
  \partial L=x\,, \tag{3.84}
\end{equation}
这与图 9 非常相像。于是，场强 $F$ 给出一个非真局域算符，它完全屏蔽伴随 Wilson 线 $\cW_1$。

\para{Wilson 线的屏蔽。} 类似地，通过取场强的乘积，我们可以构造出完全屏蔽任何整数自旋 Wilson 线
\begin{equation}
  \cW_j\,,\qquad j\in\Z \tag{3.85}
\end{equation}
的非真局域算符。更一般地，我们可以构造出把自旋相差一个整数的任意两条 Wilson 线联系起来的非真局域算符
\begin{equation}
  \cW_j\sim \cW_{j'}\,,\qquad j'-j\in\Z\,. \tag{3.86}
\end{equation}
于是未屏蔽 Wilson 线的群为
\begin{equation}
  \mathsf{D}_1=\Z_2\,. \tag{3.87}
\end{equation}
具体地说，平凡 Wilson 线与自旋-$\frac{1}{2}$ Wilson 线可以被选作构成群 $\widehat{G}^{(1)}=\mathsf{D}_1$ 的两个非平凡等价类的代表。

% [Tutor Note: p.40 式 (3.85)--(3.87) 的判别法很实用：屏蔽由“场强 $F$ 的幂”实现，而 $F$ 在自旋-1 表示中，故 $\cW_j$ 与 $\cW_{j'}$ 等价当且仅当 $j-j'\in\Z$。于是 $\mathsf{D}_1=\Z/2\Z$，代表元取 $j=0$ 与 $j=\frac{1}{2}$。核对点：$j$ 的半整数取值 (3.81) 与 $\Z_2$ 分级恰好同构，这与 $\SU(2)$ 的中心 $\Z_2$ 所对应的 1-形式对称性一致（与 3.3.2 节的 $\SO(3)$ 情形对照看更清楚）。另注：p.40 正文“the two non-trivial equivalence classes”中 “non-trivial” 疑为笔误——平凡 Wilson 线代表的正是两个类中的平凡类，只有自旋-$\frac12$ 那个类才是非平凡的；语义上应读作“这两个（互不相同的）等价类”，译文照实保留原文措辞。]

\para{电 1-形式对称性。} 取 $\mathsf{D}_1$ 的 Pontryagin 对偶，我们得到 $\SU(2)$ 纯 Yang--Mills 理论的电 1-形式对称性群
\begin{equation}
  G^{(1)}=\widehat{\mathsf{D}}_1=\Z_2\,. \tag{3.88}
\end{equation}

\para{'t Hooft 算符。} 非阿贝尔规范群 $\cG$ 的 't Hooft 算符，可以简单地由 (2.81) 附近讨论过的 $\U(1)$ 的 't Hooft 算符得到。这一点在下面的方框中说明。

\begin{notebox}{'t Hooft 算符与余特征标}
回忆 $\U(1)$ 的 't Hooft 算符 $\cH$ 具有这样的性质：在与 $\cH$ 环绕的小球面 $S^2$ 上，有
\begin{equation}
  \int_{S^2}F=m\in\Z \tag{3.89}
\end{equation}
% [原文 p.41]
其中 $F$ 是 $\U(1)$ 的场强，$m$ 是 $\cH$ 的磁荷。现在的想法是，通过让上述 $\U(1)$ 位于 $\cG$ 之内，把 $\cH$ 提升为规范群 $\cG$（假定为连通李群）的 't Hooft 算符 $\cH_\phi$。具体做法如下。

我们首先选取一个同态
\begin{equation}
  \phi : U(1)\to\cG \tag{3.90}
\end{equation}
称为\textbf{余特征标}（cocharacter）。它定义一个子群
\begin{equation}
  U(1)_\phi=\mathrm{Image}(\phi)\subseteq\cG\,. \tag{3.91}
\end{equation}
此外我们还得到一个整数
\begin{equation}
  m_\phi\in\Z \tag{3.92}
\end{equation}
描述映射 $\phi$ 绕 $\U(1)_\phi$ 的\emph{绕数}。

给定一个余特征标 $\phi$，我们就得到 $\cG$ 规范理论的一个 't Hooft 算符 $\cH_\phi$。't Hooft 算符 $\cH_\phi$ 沿时空的一个余维-3 子流形 $M_{d-3}$ 的插入
\begin{equation}
  \cH_\phi(M_{d-3})\,, \tag{3.93}
\end{equation}
是如下定义的：从时空中挖去 $M_{d-3}$，并要求路径积分中的 $\cG$ 规范场在 $M_{d-3}$ 附近具有这样的边界条件，使得在与 $M_{d-3}$ 环绕的小球面 $S^2$ 上有
\begin{equation}
  \int_{S^2}F_\phi=m_\phi \tag{3.94}
\end{equation}
其中 $F_\phi$ 是 $\U(1)_\phi$ 的场强。
\end{notebox}

在当前讨论的语境中，各种 't Hooft 算符因此由余特征标
\begin{equation}
  \phi : U(1)\to\cG=SU(2) \tag{3.95}
\end{equation}
参数化。

% [Tutor Note: p.40--41 的“余特征标”（cocharacter）即从 $\U(1)$ 到 $\cG$ 的李群同态，它给出 $\cG$ 中的一个 $\U(1)$ 子群与一个整数绕数 $m_\phi$，正是 't Hooft 算符的磁通量子化数据。核对：对 $\cG=\SU(2)$，余特征标 $\phi(z)=\mathrm{diag}(z^m,z^{-m})$ 的绕数就是 $m$，与 (3.92)、(3.94) 对应；这也说明 (3.94) 中的 $F_\phi$ 只沿 $\U(1)_\phi$ 方向取分量（即 (2.81) 的 $\U(1)$ 情形被“嵌入”到 $\cG$ 中）。]

\para{'t Hooft 算符的屏蔽。} 正如 't Hooft--Polyakov 磁单极子解那样，我们可以构造出满足性质 (3.94) 的 $\SU(2)$ 规范场的动力学解。这对任意选取的 $\phi$ 都能做到，因为唯一的要求是：由 $\phi$ 定义的 $\SU(2)$ 群流形中的圈可以缩为平凡圈。由于 $\SU(2)$ 的第一同伦群是平凡的，
\begin{equation}
  \pi_1(SU(2))=1\,, \tag{3.96}
\end{equation}
每一个这样的圈确实都是可缩的。

这些解的创生算符提供完全屏蔽余维-3 't Hooft 算符的非真余维-4 算符。于是，所有 't Hooft 算符都可以被完全屏蔽，从而
\begin{equation}
  \mathsf{D}_{d-3}=1\,. \tag{3.97}
\end{equation}

% [原文 p.42]
\para{磁 $(d-3)$-形式对称性。} 由于不存在未被屏蔽的 't Hooft 算符，也就不存在磁 $(d-3)$-形式对称性：
\begin{equation}
  G^{(d-3)}=\widehat{\mathsf{D}}_{d-3}=1\,. \tag{3.98}
\end{equation}

\para{拓扑余维-2 算符。} 上面我们推证了非平凡电 1-形式对称性 (3.88) 的存在。生成这一对称性的拓扑余维-2 算符 $U$，可理解为一个孤子型缺陷\footnote{这一术语由 [74] 引入，用于指在自身周围生成场的孤子位形的局域与延展算符。}，它在 $\SU(2)$ 理论中诱导出一个
\begin{equation}
  SO(3)=SU(2)/\Z_2 \tag{3.99}
\end{equation}
涡旋位形，而该位形无法提升为 $\SU(2)$ 涡旋位形。

更详细地说，考虑算符 $U$ 沿时空的一个余维-2 子流形 $\Sigma_{d-2}$ 的插入
\begin{equation}
  U(\Sigma_{d-2})\,, \tag{3.100}
\end{equation}
此时 $\Sigma_{d-2}$ 附近的规范场取值于 $\SO(3)$，使得在与 $\Sigma_{d-2}$ 交于一点的小圆盘 $D^2$ 上我们有
\begin{equation}
  \int_{D_2}w_2=1\in\Z/2\Z\equiv\Z_2 \tag{3.101}
\end{equation}
其中 $w_2$ 是 $\SO(3)$ 丛的\emph{第二 Stiefel--Whitney 类}（见下文）。$w_2$ 的非平凡性意味着 $\Sigma_{d-2}$ 附近的 $\SO(3)$ 规范场无法提升为 $\SU(2)$ 规范场。见图 10。

算符 $U$ 被称为 \textbf{Gukov--Witten 算符} [78,79]，它也可以被理解为某个 $\SO(3)$ 磁单极子位形的 \textbf{Dirac 弦}，而该位形无法提升为 $\SU(2)$ 磁单极子位形。作为 Dirac 弦，它可以被放在任何地方，这意味着 $U$ 是拓扑的。并且算符 $U$ 不是平凡的，因为它可以被 $\SU(2)$ 基础表示中的 Wilson 线算符探测到。

\begin{notebox}{障碍类}
考虑李群
\begin{equation}
  \cG=G/\cZ \tag{3.102}
\end{equation}
的主丛，其中 $\cZ$ 是另一个李群 $G$ 的中心 $Z(G)$ 的一个子群。$\cG$ 的主丛可以用取值于 $\cG$、定义在余维-1 位点上的转移函数来描述。这些余维-1 位点汇聚到一起，形成余维-2 的联结。考虑把这些取值于 $\cG$ 的转移函数提升为取值于 $G$ 的转移函数。在提升之前，绕一个余维-2 联结一周的转移函数之积为 1。
% [原文 p.44]
提升之后，绕一个余维-2 联结一周的转移函数之积将是 1 的一个提升，换句话说，是 $\cZ$ 中的一个元素。见图 11。由 Poincaré 对偶，这定义一个取值于 $\cZ$ 的 2-上链 $w_2$。关于 Poincaré 对偶的更多细节，见 (4.31) 附近的讨论。

$w_2$ 的上同调类被称为所研究的 $\cG$ 丛提升为 $G$ 丛的障碍类。它不依赖于上面所做的各种选择，例如提升的选取。

对
\begin{equation}
  \cG=SO(3)\,,\qquad G=SU(2)\,,\qquad \cZ=\Z_2\,, \tag{3.103}
\end{equation}
这一障碍类也被称为\textbf{第二 Stiefel--Whitney 类}。
\end{notebox}

% [原文 p.43]
\begin{figure}[htbp]
\centering
\includegraphics[width=0.36\textwidth]{fig10.png}
\figtitle{10}{生成 $\Z_2$ 电 1-形式对称性的余维-2 拓扑算符 $U$，沿 $\Sigma_{d-2}$ 插入。在与 $\Sigma_{d-2}$ 交于一点的小 2 维圆盘 $D_2$ 上，存在被开启的非平凡第二 Stiefel--Whitney 类 $w_2$，它障碍把 $\SO(3)$ 规范场提升为 $\SU(2)$ 规范场。}
\end{figure}

\begin{figure}[htbp]
\centering
\includegraphics[width=0.92\textwidth]{fig11.png}
\figtitle{11}{实线表示实现转移函数 $f_i$ 的余维-1 位点，它们沿一个余维-2 位点汇聚到一起。当绕该余维-2 位点走一圈时，人们会得到 $f_1f_2f_3$；为得到一个自洽的 $\cG$ 丛，它必须等于 $\cG$ 中的单位元。把转移函数提升为 $G$ 中的 $\tilde f_i$ 可能破坏这一自洽条件，此时一般有 $\tilde f_1\tilde f_2\tilde f_3\in\cZ$，它指定了刻画“把 $\cG$ 丛提升为 $G$ 丛”之障碍的 2-上链 $w_2$。}
\end{figure}

% [Tutor Note: 图 11 的原文图注有笔误 “which must be equal to the the identity in $\cG$”，重复了一个 “the”；译文按“等于 $\cG$ 中的单位元”处理，不改变语义。图注内容另可核对：$\tilde f_1\tilde f_2\tilde f_3$ 落在 $\cZ=\ker(G\to\cG)$ 中，因此它给出 $w_2\in H^2(X,\cZ)$ 的代表；(3.101) 中的 $w_2$ 正是它与 $D_2$ 的配对值 $1\in\Z_2$。再注：原文对同一对象混用两种写法，且四种场合并不一致——图 10 图内的标注与 p.42 散文句 “a small disk $D^2$ intersecting $\Sigma_{d-2}$ at a point” 用 $D^2$（上标），而图 10 的图注与公式 (3.101) 用 $D_2$（下标）；译文各处分别照原文保留，未作统一。]

% [原文 p.44]
\subsubsection{纯 $\SO(3)$ Yang--Mills 理论}

现在让我们考虑一个 $d$ 维规范理论，其规范群为
\begin{equation}
  \cG=SO(3)=\frac{SU(2)}{\Z_2} \tag{3.104}
\end{equation}
且不含物质场。被模掉的 $\SU(2)$ 的 $\Z_2$ 子群正是 $\SU(2)$ 的中心。

这一情形与上面讨论的纯 $\SU(2)$ Yang--Mills 理论的情形之间的差别，看起来颇为微小，因为两者唯一的区别只在于规范群中是否包含一个离散的 $\Z_2$。然而我们将看到，这一差别会关键性地改变电与磁的高阶形式对称性。

\para{Wilson 线算符。} Wilson 线算符由规范群 $\SO(3)$ 的不可约表示参数化，而这样的表示实际上比 $\SU(2)$ 的情形要少。$\SU(2)$ 的一个给定自旋-$j$ 表示 $R_j$ 只有在中心
\begin{equation}
  \Z_2\subset SU(2) \tag{3.105}
\end{equation}
在 $R_j$ 上平凡地作用时，才下降为 $\SO(3)$ 的一个表示。$\Z_2$ 中心在表示 $R_j$ 中每个矢量上作用为一个符号，该符号为 $(-1)^{2j}$。于是我们可以把中心的这一作用写成
\begin{equation}
  R_j\to(-1)^{2j}R_j\,. \tag{3.106}
\end{equation}
因此，只有 $\SU(2)$ 的整数自旋表示才是 $\SO(3)$ 的表示：
\begin{equation}
  R_j\ \text{是}\ SO(3)\ \text{表示}\iff j\in\Z\,. \tag{3.107}
\end{equation}

% [Tutor Note: p.44 式 (3.105)--(3.107) 的判据来自中心 $\Z_2=\{\pm 1\}$ 的表示论：$-1$ 在自旋-$j$ 表示上作用为 $(-1)^{2j}$，故表示下降为 $\SO(3)$ 表示当且仅当 $2j\in2\Z$，即 $j\in\Z$。核对点：$j=\frac12$ 时 $(-1)^{2j}=-1$ 非平凡，$R_{1/2}$ 不是 $\SO(3)$ 表示；$j=1$ 时 $(-1)^{2}=+1$ 平凡，$R_1$（伴随表示）是。这条判据直接决定了 $\SO(3)$ 理论中 Wilson 线的可许自旋，从而改变 $\mathsf{D}_1$ 与电 1-形式对称性（与 (3.87) 的 $\SU(2)$ 情形对照）。]
